\section{Introduction\label{sec:introduction}}


\subsection{Motivation}


\subsection{Problem Statement\label{sec:problem_statement}}

\subsection{Research Objective}

The main objective of this thesis is . 
The concept shall be based on thorough
research of . 
The objective can be
broken down into the following high-level goals:

\begin{itemize}
  \item \textbf{Obj1\label{ob:1}}: Objective
  \item \textbf{Obj2\label{ob:2}}: Objective
  \end{itemize}

\subsection{Research Questions}

The following research questions have been defined to evaluate whether an AI-assisted Process Companion, based on a \gls{rag} architecture , can effectively support crane maintenance engineers in accessing relevant information, executing maintenance procedures, and performing structured troubleshooting

\begin{itemize}
  \item \textbf{Q1\label{Q1}}: What are the main challenges and pain points that engineers face during crane troubleshooting, repair, and reporting processes?
  \item \textbf{Q2\label{Q2}}: How can large language models \acrshort{llm} assist engineers in diagnosing crane faults by retrieving, interpreting and summarizing relevant information from maintenance manuals and past repair logs?
  \item \textbf{Q3\label{Q3}}: How can Retrieval-Augmented Generation (RAG) enhance the accuracy, contextual relevance, and reliability of troubleshooting suggestions for different crane models and error codes?
  \item \textbf{Q4\label{Q4}}: How can the system provide step-by-step, process-aware guidance for troubleshooting, repair, and testing while ensuring compliance with safety and maintenance standards?
  \item \textbf{Q5\label{Q5}}: What methods and mechanisms can ensure traceability, reliability, and consistency of AI-generated outputs across multiple maintenance sessions?
  \item \textbf{Q6\label{Q6}}: How can the AI-assisted system automatically document maintenance and testing activities while maintaining organizational and regulatory compliance requirements?
  \item \textbf{Q7\label{Q7}}: How can engineers efficiently capture, store, and reuse newly observed troubleshooting experiences or field insights to continuously improve the shared knowledge base?
  \item \textbf{Q8\label{Q8}}: To what extent does the proposed LLM-based guidance system improve process efficiency, accuracy, and documentation quality compared to traditional manual methods?
  \item \textbf{Q9\label{Q9}}: What factors influence user trust, usability perception, and acceptance of AI-assisted troubleshooting and reporting systems among field engineers?
  
 \end{itemize}
These questions serve as a foundation for evaluating the role of a Process Guidance for crane maintenance engineers. The answers will help identify gaps in current methodologies and guide the development of more effective tools.


\subsection{Research Scope and Limitations}

The goal of this thesis is not to 

\clearpage
\subsection{Research Approach}

The research approach used in this thesis project is a combination of the Constructive Research
Methodology from~\cite{Lukka:2003} and the Quality Improvement Paradigm from~\cite{Wohlin:2006}.

The approach consists of the following steps that can also be repeated iteratively:

\begin{enumerate}
    \item \textbf{Find a Problem.} Find an interesting and practically relevant problem that also
        has the potential for theoretical contribution.
    \item \textbf{Understand.} Establish a baseline and obtain a practical and theoretical
        understanding of both the concrete problem and the topic area. This includes the
        elicitation of the state of practice via e.g.~interviews or observations and a
        literature review to get an overall picture of the relevant body of knowledge.
    \item \textbf{Set goals.} Define quantifiable goals that can later be used to validate the
        solution.
    \item \textbf{Innovate.} Use the obtained knowledge to construct an innovative solution
        concept and implement a proof of concept solution.
    \item \textbf{Analyse.} Analyse the solution and demonstrate that it works. Examine the scope
        of applicability and generalise. Reason how the solution can be applied to other problem
        domains or organisations. Show theoretical contributions and novelty.
\end{enumerate}


\subsection{Thesis Structure}

The remainder of this thesis is structured as follows:

\textbf{Chapter \ref{sec:foundation}} introduces the body of knowledge that represents the
foundation of this thesis. This includes relevant engineering approaches, concepts, techniques and tools.

\textbf{Chapter \ref{sec:concept}} makes use of the information from Chapter \ref{sec:foundation}
and elaborates a solution concept. This is done by summarising the state of practice,
analysing the problem in detail and exploring various solution approaches. Finally, a concrete
solution concept is presented.

\textbf{Chapter \ref{sec:realisation}} demonstrates soundness and feasibility of the elaborated
solution concept by applying the developed methods and techniques to a concrete use case from
the \gls{ELT} project. It also explains important and interesting implementation aspects.

\textbf{Chapter \ref{sec:conclusion}} summarises the results and provides an outlook on future
activities.

\cite{Apel:2013}